\documentclass[11pt,a4paper]{article}
\usepackage[margin=1in]{geometry}
\usepackage{amsmath,amssymb,amsthm}
\usepackage{algorithm}
\usepackage{algpseudocode}
\usepackage{booktabs}
\usepackage{hyperref}

\theoremstyle{definition}
\newtheorem{definition}{\textbf{Definition}}[section]
\newtheorem{theorem}{\textbf{Theorem}}[section]
\newtheorem{lemma}[theorem]{\textbf{Lemma}}
\newtheorem{remark}{\textbf{Remark}}[section]

\title{\textbf{Pipeline Parallelism: GPipe, 1F1B Scheduling, and Bubble Optimization}}
\author{\textbf{Team Scaibu} \\ \small Email: \texttt{jaydeep.9664920749@gmail.com} \\ \small \textbf{Architecture and Core Engine Team}}
\date{\textbf{October 2026}}

\begin{document}
\maketitle

\section{Definitions and Cognitive Mental Models}

\subsection{Formal Mathematical Definition}
\textbf{Definition (Pipeline Parallelism and 1F1B Scheduling):}
Pipeline Parallelism (PP) partitions a deep neural network sequentially across $P$ pipeline stages placed on distinct GPUs or cluster nodes. A training batch of size $B$ is divided into $M$ microbatches ($M \gg P$), which stream through the stages in an assembly-line fashion.

Two foundational schedules govern microbatch execution:
\begin{enumerate}
    \item \textbf{GPipe (Huang et al., 2019)}:
    Executes all $M$ forward passes sequentially, followed by all $M$ backward passes. The time spent idle while the pipeline fills (warmup) and drains (cooldown) is termed the \textbf{pipeline bubble}:
    {\boldmath
    \begin{equation}
    t_{\text{bubble}}^{\text{GPipe}} = (P - 1) \cdot (t_f + t_b)
    \end{equation}
    }
    where $t_f$ and $t_b$ are the forward and backward execution times per microbatch. The bubble fraction is:
    {\boldmath
    \begin{equation}
    F_{\text{bubble}} = \frac{P - 1}{M + P - 1}
    \end{equation}
    }
    GPipe requires storing forward activations for all $M$ microbatches on stage 0, leading to unsustainable activation memory explosion for large $M$.
    \item \textbf{1F1B (One-Forward-One-Backward) Schedule (Narayanan et al., 2021)}:
    After a warmup phase where stage $p$ runs $P - p$ forward microbatches, each stage alternates exactly one forward pass with one backward pass. As soon as a backward pass completes, the corresponding activation memory is immediately deallocated. This strictly caps the peak activation memory to $P$ microbatches regardless of total batch size $M$, enabling arbitrarily large global batches.
\end{enumerate}

\subsection{Expanded Non-Technical Definition (Intuition for Anyone)}
\begin{enumerate}
    \item \textbf{An Automobile Assembly Line}: If 4 workers build a car (Frame $\to$ Engine $\to$ Doors $\to$ Paint), worker 4 cannot start until car #1 reaches them. To keep everyone busy, cars are sent down the line one after another in microbatches.
    \item \textbf{The Bubble Problem}: At the very beginning of the day, workers 2, 3, and 4 wait idly for the first car to reach them. At the end of the day, workers 1, 2, and 3 wait idly as the last car leaves. That wasted idle time is the ``pipeline bubble''.
    \item \textbf{1F1B Memory Capping}: GPipe stores all unfinished car parts in the hallway until the whole batch finishes, filling the entire factory with clutter. 1F1B finishes one car part and immediately paints it, keeping only 4 cars in the factory at any time!
\end{enumerate}

\subsection{Child-Friendly Mental Model (Physical Metaphor)}
Imagine a bucket brigade putting out a fire:
\begin{itemize}
    \item \textbf{Bucket Delay}: Person 1 fills a bucket, hands it to Person 2, then Person 3, then Person 4. Person 4 waits 10 seconds before the first water arrives!
    \item \textbf{1F1B Balance}: Person 2 passes a full bucket forward, and immediately grabs an empty bucket passed backward. Hands are never empty, and water flows continuously!
\end{itemize}

\section{Core Mathematical Equations and Definitions}

\subsection{1F1B Pipeline Bubble Fraction}
{\boldmath
\begin{equation}
F_{\text{bubble}} = \frac{P - 1}{M + P - 1} \approx \frac{P - 1}{M} \quad (\text{for } M \gg P)
\end{equation}
}
\begin{itemize}
    \item \textbf{Formal Definition}: Ratio of idle processor cycles to total cluster capacity.
    \item \textbf{What It Means}: Bubble overhead shrinks toward zero as microbatch count $M$ increases relative to stages $P$.
    \item \textbf{Why It Is Important}: Dictates the minimum batch size required to achieve $>90\%$ cluster utilization.
\end{itemize}

\subsection{Peak Stored Activation Memory}
{\boldmath
\begin{align}
\mathcal{M}_{\text{act}}^{\text{GPipe}} &= M \cdot a_{\text{micro}} \\
\mathcal{M}_{\text{act}}^{\text{1F1B}} &= P \cdot a_{\text{micro}}
\end{align}
}
\begin{itemize}
    \item \textbf{Formal Definition}: Maximum forward activation cache size where $a_{\text{micro}}$ is activation footprint per microbatch.
    \item \textbf{What It Means}: 1F1B decouples memory from batch size $M$, binding it strictly to stage count $P$.
    \item \textbf{Why It Is Important}: Primary architectural reason 1F1B is the universal standard in Megatron-LM.
\end{itemize}

\section{Rigorous Mathematical Analysis Checklist}

\begin{itemize}
    \item \textbf{Notation}: $P$ (pipeline stages), $M$ (microbatches), $t_f, t_b$ (forward/backward step times).
    \item \textbf{Epistemic Status}: Huang et al. (GPipe, 2019); Narayanan et al. (Megatron-LM, 2021). Universal multi-node scaling primitive.
    \item \textbf{Dependency Graph}: Microbatch $m \to \text{Stage } 0 \to \dots \to \text{Stage } P-1 \to \text{Loss} \to \text{Backward } P-1 \to \dots \to \text{Stage } 0$.
    \item \textbf{Dimensions}: Stage count $P \in [4, 64]$, microbatches $M \ge 4P$.
    \item \textbf{Regimes}: Cross-node cluster scaling combined with Tensor Parallelism and Data Parallelism (3D Parallelism).
    \item \textbf{Invariants}: Point-to-point activations communicated between stage $p$ and $p+1$ only; zero full-cluster All-Reduces.
    \item \textbf{Isomorphisms}: Isomorphic to classic hardware instruction pipelining with hazard stalls.
\end{itemize}

\section{Theoretical Theorems and Formal Proofs}

\begin{theorem}[Bubble Overhead Lower Bound Under 1F1B Scheduling]
Let a model be partitioned across $P$ stages with balanced forward time $t_f$ and backward time $t_b = 2 t_f$. For $M \ge P$ microbatches, the total execution time $T_{\text{total}}$ under 1F1B is strictly bounded by:
{\boldmath
\begin{equation}
T_{\text{total}} = M(t_f + t_b) + (P - 1)(t_f + t_b)
\end{equation}
}
and the bubble fraction is identically $F_{\text{bubble}} = \frac{P - 1}{M + P - 1}$.
\end{theorem}

\begin{proof}
Let the timeline be measured in unit cycles where forward pass takes $t_f$ and backward pass takes $t_b$.
1. Stage $P - 1$ (the final stage) cannot begin processing microbatch 0 until it traverses stages $0, 1, \dots, P - 2$.
The warmup delay on Stage $P - 1$ is:
{\boldmath
\begin{equation}
T_{\text{warmup}} = (P - 1) t_f
\end{equation}
}
2. Once Stage $P - 1$ begins, it continuously processes all $M$ forward and backward microbatches with zero internal stalls under 1F1B scheduling, taking:
{\boldmath
\begin{equation}
T_{\text{active}} = M(t_f + t_b)
\end{equation}
}
3. After Stage $P - 1$ completes the final backward pass of microbatch $M - 1$, the backward activations must travel back through stages $P - 2, \dots, 0$. The cooldown delay on Stage 0 is:
{\boldmath
\begin{equation}
T_{\text{cooldown}} = (P - 1) t_b
\end{equation}
}
The total makespan is the sum of warmup delay, active processing, and cooldown delay:
{\boldmath
\begin{equation}
T_{\text{total}} = T_{\text{warmup}} + T_{\text{active}} + T_{\text{cooldown}} = (P - 1) t_f + M(t_f + t_b) + (P - 1) t_b = M(t_f + t_b) + (P - 1)(t_f + t_b)
\end{equation}
}
The ideal computation time without idle bubbles is $T_{\text{ideal}} = M(t_f + t_b)$.
The bubble fraction is:
{\boldmath
\begin{equation}
F_{\text{bubble}} = \frac{T_{\text{total}} - T_{\text{ideal}}}{T_{\text{total}}} = \frac{(P - 1)(t_f + t_b)}{M(t_f + t_b) + (P - 1)(t_f + t_b)} = \frac{P - 1}{M + P - 1}
\end{equation}
}
This completes the formal proof.
\end{proof}

\begin{theorem}[Bubble Reduction in Interleaved 1F1B Schedules]
If each device is assigned $v \ge 2$ virtual non-adjacent stages (interleaved 1F1B), the bubble fraction is reduced by factor $v$:
{\boldmath
\begin{equation}
F_{\text{bubble}}^{\text{interleaved}} \approx \frac{1}{v} \cdot \frac{P - 1}{M}
\end{equation}
}
\end{theorem}

\begin{proof}
In interleaved 1F1B, each stage contains $1/v$ of a physical stage's layers, so the forward microbatch duration scales down to $t_f' = t_f / v$ and backward duration to $t_b' = t_b / v$.
The warmup idle time for a microbatch to traverse $P - 1$ virtual stages is:
{\boldmath
\begin{equation}
T_{\text{bubble}}^{\text{interleaved}} = (P - 1) (t_f' + t_b') = \frac{(P - 1)(t_f + t_b)}{v}
\end{equation}
}
Dividing by ideal execution time $M(t_f + t_b)$:
{\boldmath
\begin{equation}
F_{\text{bubble}} = \frac{1}{v} \cdot \frac{P - 1}{M}
\end{equation}
}
For $v = 2$, the pipeline bubble is halved, reducing wasted idle time by 50\% at the expense of doubling point-to-point network messages.
\end{proof}

\section{Foundational Architectural Questions and Epistemic Meta-Questions}

\subsection{Architectural Question 1: Stage Workload Balancing}
\textbf{Question:} Why does partitioning an 80-layer model into 8 equal 10-layer stages cause massive pipeline bubbles?\\
\textbf{Answer:} Stage 0 contains the heavy input token embedding layer ($V \times d$), while stage 7 contains the vocabulary output projection and cross-entropy loss computation ($d \times V$). If layers are divided equally (10 layers per stage), Stage 0 and Stage 7 take $30\%\text{--}50\%$ longer to compute than interior stages 1–6, stalling the entire pipeline behind the slowest stage. Stages must be balanced by measured execution time (e.g. 8 layers on stage 0, 11 layers on interior stages, 8 layers on stage 7).

\subsection{Architectural Question 2: Point-to-Point Network Communication}
\textbf{Question:} Why is Pipeline Parallelism vastly more network-efficient than Tensor Parallelism across nodes?\\
\textbf{Answer:} Tensor Parallelism requires all GPUs to communicate simultaneously via high-overhead All-Reduces. Pipeline Parallelism requires communication only between adjacent stages (GPU $p$ sends to GPU $p+1$), transferring only intermediate activation boundary tensors ($B \times L \times d$) over standard point-to-point network links.

\section{Algorithmic Implementation Specification}

\begin{algorithm}[H]
\caption{1F1B Pipeline Bubble and Memory Evaluation}
\begin{algorithmic}[1]
\Require Stage count $P$, microbatches $M$, forward duration $t_f$, backward duration $t_b$, schedule
\Ensure Bubble fraction $F$, total duration $T$, peak activation memory $A_{\max}$, ideal duration $T_{\text{id}}$
\State $T_{\text{id}} \gets M \cdot (t_f + t_b)$
\If{$\text{schedule} == \text{"interleaved"}$}
    \State $t_{\text{bubble}} \gets 0.5 \cdot (P - 1) \cdot (t_f + t_b)$
    \State $A_{\max} \gets P$
\ElsIf{$\text{schedule} == \text{"gpipe"}$}
    \State $t_{\text{bubble}} \gets (P - 1) \cdot (t_f + t_b)$
    \State $A_{\max} \gets M$
\Else
    \State $t_{\text{bubble}} \gets (P - 1) \cdot (t_f + t_b)$
    \State $A_{\max} \gets P$
\EndIf
\State $T \gets T_{\text{id}} + t_{\text{bubble}}$
\State $F \gets t_{\text{bubble}} / T$
\State \Return $(F, T, A_{\max}, T_{\text{id}})$
\end{algorithmic}
\end{algorithm}

\section{Big-$\Theta$ Complexity Analysis}

\begin{itemize}
    \item \textbf{Pipeline Schedule Calculation Time}: $\Theta(1)$ constant mathematical evaluations.
    \item \textbf{Peak Memory Footprint}: $\Theta(P)$ microbatch activations under 1F1B.
    \item \textbf{Bubble Latency}: $\mathcal{O}(P / M)$ asymptotic decay.
    \item \textbf{Communication Complexity}: $\Theta(M \cdot B \cdot L \cdot d)$ point-to-point floats per stage boundary.
\end{itemize}

\section{Single Canonical Repository Reference}

The complete deterministic scalar implementation, verification suite, and unit test benchmarks are published at:
\begin{center}
\url{https://github.com/Chief-Strategist-J/llm-obs-infra}
\end{center}

\end{document}
